\documentclass[12pt, a4paper]{article}
\usepackage{fontspec}
\usepackage{geometry}
\usepackage{hyperref}
\usepackage{booktabs}
\usepackage{enumitem}
\usepackage{xcolor}
\usepackage{listings}
\usepackage{fancyhdr}
\usepackage{titlesec}
\usepackage{wasysym}
\setmainfont{DejaVu Serif}
\setsansfont{DejaVu Sans}
\setmonofont{DejaVu Sans Mono}

\geometry{margin=2.5cm}

\definecolor{codegray}{RGB}{245,245,245}
\definecolor{wujia}{RGB}{30,130,80}

\lstset{
  backgroundcolor=\color{codegray},
  basicstyle=\ttfamily\small,
  breaklines=true,
  frame=single,
  rulecolor=\color{gray},
}

\pagestyle{fancy}
\fancyhf{}
\rhead{WujiaTea ERP}
\lhead{Checklist triển khai server}
\rfoot{Trang \thepage}

\titleformat{\section}{\large\bfseries\color{wujia}}{}{0em}{}[\titlerule]

\begin{document}

\begin{center}
  {\LARGE\bfseries WujiaTea ERP — Checklist triển khai}\\[6pt]
  {\small\color{gray} Cập nhật: 30/04/2026}
\end{center}

\vspace{1em}

% ── Bước 1 ─────────────────────────────────────────────────────────
\section{Bước 1 — Code \& Git}

\begin{itemize}[label=\CheckedBox]
  \item Tạo GitHub repo \texttt{huyban2003/WujiaTea-ERP-internal}
  \item Push toàn bộ code lên GitHub
\end{itemize}
\begin{itemize}[label=\Square]
  \item Đổi repo về \textbf{Private} sau khi xong
\end{itemize}

% ── Bước 2 ─────────────────────────────────────────────────────────
\section{Bước 2 — Setup Windows Server (qua RDP)}

\textit{Kết nối: Remmina $\to$ RDP $\to$ \texttt{113.161.187.126:5761} $\to$ user \texttt{dev}}

\begin{itemize}[label=\Square]
  \item Cài \textbf{Git for Windows}: \url{https://git-scm.com/download/win}
  \item Cài \textbf{Docker Desktop}: \url{https://www.docker.com/products/docker-desktop}
    \begin{itemize}
      \item Restart máy sau khi cài
      \item Bật \textit{``Use WSL 2 based engine''} trong Settings
    \end{itemize}
  \item Clone repo về server:
\begin{lstlisting}
git clone https://github.com/huyban2003/WujiaTea-ERP-internal.git C:\wujia-tea
\end{lstlisting}
  \item Bật OpenSSH Server (chạy PowerShell Admin):
\begin{lstlisting}
Add-WindowsCapability -Online -Name OpenSSH.Server~~~~0.0.1.0
Start-Service sshd
Set-Service -Name sshd -StartupType Automatic
New-NetFirewallRule -Name "OpenSSH" -DisplayName "OpenSSH Server" `
  -Enabled True -Direction Inbound -Protocol TCP -Action Allow -LocalPort 22
\end{lstlisting}
\end{itemize}

% ── Bước 3 ─────────────────────────────────────────────────────────
\section{Bước 3 — Chạy Odoo bằng Docker}

\begin{itemize}[label=\Square]
  \item Mở PowerShell tại \texttt{C:\textbackslash wujia-tea} rồi chạy:
\begin{lstlisting}
docker compose up -d
\end{lstlisting}
  \item Truy cập \texttt{http://113.161.187.126:8019} để kiểm tra
  \item Tạo database \texttt{wujia\_tea\_19} qua Odoo web interface
\end{itemize}

% ── Bước 4 ─────────────────────────────────────────────────────────
\section{Bước 4 — Domain \& SSL}

\begin{itemize}[label=\Square]
  \item Mua/dùng domain (ví dụ: \texttt{erp.wujiatea.com})
  \item Trỏ DNS \texttt{A record} $\to$ \texttt{113.161.187.126}
  \item Mở port 80 và 443 trên FortiGate (NAT vào server)
  \item Sửa \texttt{config/nginx.conf}: thay \texttt{wujiatea.com} bằng domain thật
  \item Lấy SSL:
\begin{lstlisting}
docker compose run --rm certbot certonly ^
  --webroot -w /var/www/certbot -d erp.wujiatea.com
\end{lstlisting}
  \item Restart nginx: \texttt{docker compose restart nginx}
  \item Kiểm tra HTTPS hoạt động
\end{itemize}

% ── Bước 5 ─────────────────────────────────────────────────────────
\section{Bước 5 — CI/CD GitHub Actions}

\begin{itemize}[label=\Square]
  \item Vào GitHub repo $\to$ \textbf{Settings $\to$ Secrets $\to$ Actions}
  \item Thêm 3 secrets:
    \begin{itemize}
      \item \texttt{SERVER\_HOST} = \texttt{113.161.187.126}
      \item \texttt{SERVER\_USER} = \texttt{dev}
      \item \texttt{SERVER\_PASS} = \texttt{Dev@2026}
    \end{itemize}
  \item Push 1 commit lên \texttt{main} $\to$ vào tab \textbf{Actions} kiểm tra pipeline
\end{itemize}

% ── Bước 6 ─────────────────────────────────────────────────────────
\section{Bước 6 — Bảo mật sau khi xong}

\begin{itemize}[label=\Square]
  \item Đổi GitHub repo về \textbf{Private}
  \item Đổi password \texttt{Dev@2026} trên server
  \item Xoá/rotate GitHub token (xem \texttt{docs/CREDENTIALS.md})
  \item Cấu hình backup database định kỳ
\end{itemize}

\end{document}
